\clearpage\corrigesChapitre{7}{Changements d’état et énergie}
\corrigesSouples
\begin{blocLivre}{cor-7-1}\begin{corrige}{1}{Nommer les changements}
\begin{questions}
\item Liquéfaction de la vapeur d’eau sur la vitre ; fusion du glaçon.
\item Non, les entités restent des molécules de \ce{H2O}.
\end{questions}
\end{corrige}\end{blocLivre}
\begin{blocLivre}{cor-7-2}\begin{corrige}{2}{Écrire symboliquement}
\begin{questions}
\item $\ce{H2O(l) -> H2O(g)}$.
\item La vaporisation est endothermique : l’eau reçoit de l’énergie.
\end{questions}
\end{corrige}\end{blocLivre}
\begin{blocLivre}{cor-7-3}\begin{corrige}{3}{Calculer une énergie de fusion}
\begin{questions}
\item $m=\qty{0,0500}{kg}$.
\item $E=mL_f=0,0500\times3,34\times10^5=\qty{1,67e4}{J}$.
\end{questions}
\end{corrige}\end{blocLivre}
\begin{blocLivre}{cor-7-4}\begin{corrige}{4}{Dissoudre ou fondre ?}
\begin{questions}
\item Le sucre se dissout ; la cire fond.
\item La dissolution fait intervenir l’eau comme solvant.
\end{questions}
\end{corrige}\end{blocLivre}
\begin{blocLivre}{cor-7-5}\begin{corrige}{5}{Lire une courbe thermique}
\schema[corrige]{palier}
\begin{questions}
\item Le palier à 0 °C correspond à la fusion dans l’exemple de l’eau.
\item L’énergie reçue permet le changement d’organisation du solide en liquide, au lieu d’élever la température du corps pur.
\end{questions}
\end{corrige}\end{blocLivre}
\begin{blocLivre}{cor-7-6}\begin{corrige}{6}{Une masse à solidifier}
\begin{questions}
\item $m=E/L_f=66,8/334=\qty{0,200}{kg}$.
\item Le milieu extérieur reçoit l’énergie cédée par l’eau ; la solidification est exothermique.
\end{questions}
\end{corrige}\end{blocLivre}
\begin{blocLivre}{cor-7-7}\begin{corrige}{7}{Deux masses de glace}
\begin{questions}
\item $E_2/E_1=m_2/m_1=100/40=2,5$.
\item Oui : c’est le même changement d’état de la même espèce dans les mêmes conditions.
\end{questions}
\end{corrige}\end{blocLivre}
\begin{blocLivre}{cor-7-8}\begin{corrige}{8}{Une évaporation qui refroidit}
\begin{questions}
\item L’eau reçoit de l’énergie, notamment de la peau, pour se vaporiser.
\item La peau cède de l’énergie thermique ; sa température peut diminuer si cet apport n’est pas compensé.
\end{questions}
\end{corrige}\end{blocLivre}
\begin{blocLivre}{cor-7-9}\begin{corrige}{9}{Une puissance utile fournie}
\begin{questions}
\item $E=0,200\times3,34\times10^5=\qty{6,68e4}{J}$.
\item $t=66800/100=\qty{668}{s}$, soit environ 11,1 min. Le débit d’énergie utile est fourni dans l’énoncé.
\end{questions}
\end{corrige}\end{blocLivre}
\begin{blocLivre}{cor-7-10}\begin{corrige}{10}{Masse conservée, volume changé}
\begin{questions}
\item $V_l=\qty{100}{cm^3}$ ; $V_s=100/0,92\simeq\qty{109}{cm^3}$.
\item Le volume augmente d’environ 9 cm$^3$ ; sans espace libre, la solidification peut exercer des contraintes importantes.
\end{questions}
\end{corrige}\end{blocLivre}
\begin{blocLivre}{cor-7-11}\begin{corrige}{11}{Une étape oubliée}
\begin{questions}
\item Cette relation ne calcule que l’énergie de fusion à la température de fusion.
\item Chauffer la glace jusqu’à 0 °C ; la faire fondre ; chauffer l’eau liquide jusqu’à 20 °C.
\end{questions}
\end{corrige}\end{blocLivre}
\begin{blocLivre}{cor-7-12}\begin{corrige}{12}{Comparer deux matériaux}
\demarche{Comparer les énergies pour une même masse, puis les températures d’utilisation.}
\begin{questions}
\item B : son énergie échangée par kilogramme transformé est supérieure.
\item Il faut connaître sa température de fusion et vérifier qu’elle convient à la plage de fonctionnement.
\end{questions}
\retiens{Une grande énergie massique ne suffit pas : la température du changement d’état doit aussi convenir.}
\end{corrige}\end{blocLivre}

